% frmcs_appsession_fr.tex  — Draft v2.9
% XeLaTeX — ETSI FRMCS Concept de Session Applicative (Français)
\documentclass[11pt,a4paper]{article}

\usepackage{fontspec}
\setmainfont[SmallCapsFont = Avenir LT Std 35 Light]{Avenir LT Std 35 Light}
\setsansfont{Avenir LT Std 45 Book}

\usepackage{geometry}\geometry{a4paper,margin=2.5cm}
\usepackage{polyglossia}\setdefaultlanguage{french}
\usepackage{microtype,booktabs,longtable,array,listings,hyperref,titlesec,enumitem,soul,parskip}
\usepackage[table]{xcolor}
\usepackage{pifont}
\newcommand{\cmark}{\ding{51}}

\definecolor{codebg}{RGB}{245,245,245}
\definecolor{etsiblue}{RGB}{0,70,127}
\definecolor{noteblue}{RGB}{230,240,255}

\lstset{backgroundcolor=\color{codebg},basicstyle=\small\ttfamily,breaklines=true,
  frame=single,framerule=0.4pt,rulecolor=\color{gray!40},
  xleftmargin=6pt,xrightmargin=6pt,aboveskip=8pt,belowskip=4pt,
  showstringspaces=false,keepspaces=true}

\titleformat{\section}{\large\bfseries\sffamily\color{etsiblue}}{{\thesection}}{1em}{}
\titleformat{\subsection}{\normalsize\bfseries\sffamily\color{etsiblue}}{{\thesubsection}}{1em}{}
\titleformat{\subsubsection}{\normalsize\bfseries\sffamily}{{\thesubsubsection}}{1em}{}

\newenvironment{notebox}{%
  \begin{center}\begin{tabular}{|p{0.93\textwidth}|}\hline\rowcolor{noteblue}
}{\\\hline\end{tabular}\end{center}}

% Interface name shorthand
\newcommand{\OBAPP}{OB\textsubscript{APP}}
\newcommand{\TSAPP}{TS\textsubscript{APP}}
% Arrow substitute: robust in text and table contexts
\newcommand{\arr}{$\to$}

\newcommand{\svcref}{\texttt{svcSessionRef}}
\newcommand{\appcall}{\texttt{appSessionId\_call}}
\newcommand{\appterm}{\texttt{appSessionId\_term}}

\hypersetup{colorlinks=true,linkcolor=etsiblue,urlcolor=etsiblue,
  pdfauthor={R2DATO D25.3},pdftitle={ETSI FRMCS -- Concept de Session Applicative}}

\begin{document}

\begin{titlepage}
  \vspace*{3cm}{\sffamily\color{etsiblue}
    {\Huge\bfseries ETSI FRMCS\\[0.4em]Concept de Session Applicative}\\[1.5em]
    {\Large Proposition de modifications ETSI TS 103\,764, TS 103\,765-3/4/5}\\[2em]
    {\large Contribution de travail -- Draft v2.9}\\[0.5em]
    {\normalsize R2DATO D25.3 -- Concept Paper FRMCS}\\[4em]
    \hrule\vspace{0.5em}
    {\small\textbf{Documents cibles :} ETSI TS 103\,764 (RT-0095),
      TS 103\,765-3 (RT-0097), TS 103\,765-4 (RT-0099), TS 103\,765-5 (RT-00104)\\[0.3em]
      \textbf{Documents inchangés :} toutes les spécifications UIC (FFFIS-7950,
      FIS-7970, FRS FU-7120, AT-7800, TOBA FRS-7510, EIRENE FRS, EIRENE SRS)}}
\end{titlepage}

\tableofcontents\clearpage

% ── 1. PRINCIPE ARCHITECTURAL ───────────────────────────────────────────────
\section{Principe architectural}

\subsection{Reformulation de la notion de session}

Dans les TS ETSI, \textbf{« session »} désignait le tunnel MCData IPcon établi par la
GW (ce que la présente proposition appelle \textbf{session de service}). Redéfinition :

\begin{itemize}[leftmargin=2em]
  \item \textbf{« session »} : session applicative -- communication entre deux
    applications FRMCS, composée d'une ou plusieurs sessions de service.
  \item \textbf{« session de service »} : bearer individuel. La première est
    nécessairement MCx (MCData IPcon) -- inchangé. Les additionnelles peuvent être
    tout type MCx ou futurs services critiques ferroviaires (R2DATO D25.3 Part B).
\end{itemize}

Toutes les occurrences internes « session » \arr\ « session de service » dans les TS.
Identifiants applicatifs (\texttt{sessionId}, \texttt{POST /sessions},
\texttt{openSession*}, \texttt{closeSession*}, \texttt{incomingSessionNotif*},
\texttt{open session}) : \textbf{inchangés}.

\begin{notebox}
  \textbf{Ce qui change pour les applications : rien.} Les interfaces \OBAPP\ et
  \TSAPP\ sont strictement inchangées. Les specs UIC utilisent déjà « session » au
  sens applicatif -- leur terminologie est désormais naturellement alignée.
\end{notebox}

\subsection{Pourquoi l'ASAS plutôt que le champ AnyExt du SIP INVITE}

RT-0097 §6.2.2.3 porte déjà des données applicatives dans \texttt{AnyExt} (format
« not specified »). Techniquement faisable pour MCx, mais structurellement exclu :
un futur service critique non-MCx n'a pas de SIP INVITE. L'ASAS agnostique est le
seul mécanisme survivant à la séparation des strates (R2DATO D25.3 Part B).

\subsection{Résolution d'identifiants applicatifs (AIR)}

La fonction AIR résout le \texttt{remoteId} en \textbf{identifiant de service}
(MCx ou futur service critique).
\textbf{Règles pour la première édition -- entièrement spécifiées :}

\begin{enumerate}[leftmargin=2em]
  \item Si \texttt{remoteId} se termine par \texttt{.ertms} (format ETCS
    \texttt{id.ty.cc.ertms} défini dans UIC AT-7800~§5.3, référençant
    UNISIG SUBSET-037) : complétion par le domaine courant
    \arr\ \texttt{id.ty.cc.ertms@current\_domain}.
    Conforme aux specs UIC actuelles. Fonctionnel mais non-optimal en zone de
    chevauchement de domaines.
  \item Sinon : pass-through FA MCx -- comportement identique à l'historique.
\end{enumerate}

\begin{notebox}
  La fonction AIR est conçue comme un \textbf{point d'extension}. Les éditions
  futures pourront introduire une résolution DNS ou un mapping FROP plus précis
  sans impact sur \OBAPP, \TSAPP, SIP ni les specs UIC, en lien avec les travaux
  ERJU sur l'adressage ETCS dans FRMCS (UIC EIRENE~SRS, UNISIG SUBSET-037).
\end{notebox}

\subsection{Mécanisme d'initialisation}

Une fois la première session de service établie (SIP + GRE inchangés), chaque GW
expose un \textbf{ASAS} accessible sur son adresse IP post-NAT (SD\_PUB\_IP,
cf.\ §\ref{sec:nat}). La GW appelante pilote :

\begin{enumerate}[leftmargin=2em]
  \item \textbf{Initialisation} : POST à l'ASAS ; reçoit \appterm, clé Ed25519,
    liste d'alias, politique MSP, période keepalive.
  \item \textbf{Sessions de service additionnelles} : établies selon liste d'alias +
    FROP ; chacune déclarée à l'ASAS avec \texttt{sessionToken} signé, permettant
    l'authentification mutuelle indépendamment du protocole de signalisation
    sous-jacent.
  \item \textbf{Keepalive intelligent} : envoyé uniquement si silence applicatif
    $\geq$ T\_KA sur un \svcref\ donné (cf.\ §\ref{sec:ka}).
  \item \textbf{Mise à jour alias} : PATCH immédiat de la GW destinatrice sur tout
    nouvel enregistrement d'alias ou d'identifiant de service. Sans configuration.
  \item \textbf{Libération} : DELETE à l'ASAS + libération de chaque session de
    service selon son propre protocole de signalisation.
\end{enumerate}

\subsection{Cycle de vie de la session applicative}

Vivante tant qu'au moins une session de service est \texttt{ACTIVE} ou
\texttt{DEGRADED}. Timer de recovery : INTEGER(1..180)\,s (FROP) ;
$\leq 5$\,s pour ETCS/ATP (AT-7800 Annex A.6).
À l'expiration du timer, la \textbf{session applicative} est libérée.

% ── 2. CONCEPTS CLÉS ────────────────────────────────────────────────────────
\section{Concepts clés}

\subsection{Terminologie}

\begin{longtable}{>{\bfseries}p{3.8cm} p{10.5cm}}
  \toprule Terme & Définition \\ \midrule \endhead
  session & Communication applicative entre deux applications FRMCS, composée
    d'une ou plusieurs sessions de service. Identifiée par \texttt{sessionId}
    (\OBAPP/\TSAPP), en interne par (\appcall, \appterm). \\[0.3em]
  session de service & Bearer individuel, transparent aux applications.
    La première est nécessairement MCx (MCData IPcon). Les additionnelles
    peuvent être MCx ou futurs services critiques. Identifiée par \svcref. \\[0.3em]
  open session & Session non libérée (inchangé). \\[0.3em]
  AIR & Fonction interne résolvant \texttt{remoteId} en \textbf{identifiant de
    service} (MCx ou futur service critique). Règles v1 entièrement spécifiées :
    \texttt{*.ertms} \arr\ \texttt{*.ertms@current\_domain}
    (UIC AT-7800~§5.3, réf.\ UNISIG SUBSET-037) ;
    sinon pass-through FA MCx. Point d'extension pour une résolution future
    (DNS/FROP, en lien avec les travaux ERJU sur l'adressage ETCS --
    UIC EIRENE~SRS, UNISIG SUBSET-037). Non visible des applications. \\[0.3em]
  ASM & Fonction interne gérant le cycle de vie des sessions et sessions de
    service, pilotant les interactions ASAS et la politique MSP. \\[0.3em]
  ASAS & Serveur d'API REST exposé par chaque GW sur son IP post-NAT (SD\_PUB\_IP).
    La sécurité repose sur l'authentification par message Ed25519 ; il n'existe pas
    d'exigence normative de confidentialité des échanges ASAS entre GW. Le transport
    HTTP PEUT être protégé par TLS (recommandé). Ancrage sur la première session de
    service MCx -- pas de PKI inter-GW requise. \\[0.3em]
  MSP & Multi Service Path : mécanisme inter-GW gérant plusieurs sessions de service
    parallèles au sein d'une même session applicative. Distinct du Multipath transport
    (RT-00104). Politique (\texttt{mspPolicy}) négociée à l'init. ASAS :
    DUPLICATION, RECOVERY ou LOADBALANCE. \\[0.3em]
  GW appelante & Initie \texttt{POST /sessions} via \OBAPP.
    Dérive l'IP ASAS de l'IP source SIP post-NAT. \\[0.3em]
  GW destinatrice & Expose l'ASAS. Construit la liste d'alias ; DOIT envoyer
    immédiatement \texttt{PATCH /aliases} sur tout nouveau enregistrement ;
    notifie l'application via \TSAPP\ (\texttt{incomingSessionNotif*}). \\[0.3em]
  liste d'alias & Liste dynamique d'identifiants de service + types + domaines.
    Mise à jour immédiate via \texttt{PATCH /aliases} -- sans configuration FROP. \\[0.3em]
  identifiant de service & Identifiant pour établir une session de service, quel
    que soit le type. FA ou MC Service User ID pour MCx ; tout identifiant
    associable à l'utilisateur pour les futurs services critiques. \\[0.3em]
  timer de recovery & INTEGER(1..180)\,s (FROP) ; $\leq 5$\,s pour ETCS/ATP
    (AT-7800 Annex A.6). La session applicative est libérée à l'expiration. \\[0.3em]
  T\_KA & Période keepalive : INTEGER(1..10)\,s (application profile).
    Règle : T\_KA $\leq$ timer\_recovery\,/\,2 ($\leq 2,5$\,s pour ETCS/ATP).
    $N = 2$ échecs consécutifs (fixe) \arr\ \texttt{DEGRADED}. \\
  \bottomrule
\end{longtable}

\subsection{Identifiants de session}

\begin{itemize}[leftmargin=2em]
  \item \appcall\ (UUID v4) : GW appelante, stable pour toute la durée de la session.
  \item \appterm\ (UUID v4) : GW destinatrice à l'init. ASAS. Analogue aux TEIDs GTP.
  \item \svcref\ (UUID v4) : \textbf{Stable et immuable}, y compris lors du recovery.
  \item \texttt{svcSessionId} (natif) : SIP Call-ID pour MCx. Corrélateur ASAS.
    Peut changer lors d'un recovery.
\end{itemize}

\subsection{Modèle d'authentification}
\label{sec:trust}

\textbf{La première session de service MCx est un service sûr déjà en charge
d'identifier les parties} -- ancrage d'authentification suffisant. Il n'existe pas
d'exigence normative de confidentialité des échanges ASAS entre GW : l'ASAS opère
dans un contexte inter-GW ferroviaire où la protection des métadonnées de session
n'est pas un prérequis normatif.

\subsubsection*{Niveau 1 -- Transport HTTP (TLS recommandé)}

Le transport HTTP de l'ASAS PEUT être protégé par TLS~1.3 (recommandé pour
l'intégrité du canal). Certificat serveur \textbf{auto-signé} possible -- aucune PKI
inter-opérateur. TLS n'est pas normatif à ce niveau ; l'authentification normative
est assurée par les niveaux~2 et~3. L'option mTLS est possible sur accord bilatéral
FSNNI OOB (FROP).

\begin{notebox}
  \textbf{Note FSNNI multi-hop} : le topology hiding ne compromet pas
  l'authentification -- le \texttt{sessionToken} Ed25519 protège l'intégrité par
  message indépendamment du chemin réseau. Seule contrainte : routabilité ASAS sur
  SD\_PUB\_IP (§\ref{sec:nat}).
\end{notebox}

\subsubsection*{Niveau 2 -- Ancrage d'authentification (première session MCx)}

Le MC Service Server a authentifié la GW appelante avant d'accepter la session.
Échange Ed25519 sur ce canal = bootstrapping sûr : toute clé reçue à l'initialisation
ASAS est authentique par construction.

\subsubsection*{Niveau 3 -- Authentification par message (sessionToken Ed25519)}

\begin{lstlisting}[language={}]
sessionToken = Sign(sk_call,
  appSessionId_call || 0x00 || appSessionId_term || 0x00 ||
  svcSessionRef     || 0x00 || svcSessRequestTime || 0x00 ||
  svcSessAcceptTime)
// Recovery : ajouter 0x00 || replSvcSessionRef
ackToken = Sign(sk_term, meme entree)
\end{lstlisting}

Sérialisation : concaténation UTF-8 avec séparateurs \texttt{0x00} -- déterministe,
sans dépendance externe.

\subsection{Keepalive intelligent}
\label{sec:ka}

\begin{itemize}[leftmargin=2em]
  \item \textbf{T\_KA} : INTEGER(1..10)\,s (application profile) ;
    $\leq 2,5$\,s pour ETCS/ATP.
  \item \textbf{$N = 2$} échecs consécutifs (fixe) \arr\ \texttt{DEGRADED}.
  \item Recovery : nouvelle session avec même \svcref\ + token signé.
\end{itemize}

\subsection{Adressage ASAS et modèle NAT}
\label{sec:nat}

\textbf{Règle d'adressage ASAS (résolue)} : IP source post-NAT du signalement SIP
(SD\_PUB\_IP, RT-0097 §6.2.2.4). Aucune extraction de header SIP. Port FROP.

% ── 3. RÈGLE DE REMPLACEMENT ────────────────────────────────────────────────
\section{Règle de remplacement terminologique dans les TS ETSI}

\begin{longtable}{p{5.5cm} p{4cm} p{4.5cm}}
  \toprule Contexte & Terme actuel & Terme proposé \\ \midrule \endhead
  Bearer MCx interne & session & \textbf{session de service} \\
  Établissement MCx interne & session establishment & \textbf{service session} establishment \\
  Libération MCx interne & session release & \textbf{service session} release \\
  Multi-chemin inter-GW (service sessions) & Multipath session & \textbf{Multi Service Path (MSP)} \\
  Multipath transport (RT-00104) & Multipath & \textbf{Multipath} (inchangé) \\
  Identifiant \OBAPP/\TSAPP & \texttt{sessionId} & \textbf{inchangé} \\
  Demande application & \texttt{POST /sessions}, \texttt{openSession*} & \textbf{inchangé} \\
  Notification entrante & \texttt{incomingSessionNotif*} & \textbf{inchangé} \\
  Libération application & \texttt{closeSession*}, \texttt{DELETE /sessions} & \textbf{inchangé} \\
  État applicatif & \texttt{open session} & \textbf{inchangé} \\
  \bottomrule
\end{longtable}

% ── 4. MODIFICATIONS ETSI ───────────────────────────────────────────────────
\section{Modifications ETSI proposées}

\subsection{ETSI TS 103\,764 (RT-0095) -- Annexe normative X}

\paragraph{X.1 -- Périmètre}
API inter-GW agnostique au type de session de service. Non visible des applications.
Ne modifie aucune spec UIC ni protocole de service. Absorbe \texttt{FUTURE\_CRITICAL\_SERVICE}
sans modification de la présente annexe.

\paragraph{X.2 -- Transport et authentification}

\textbf{Transport :} HTTP/2 + JSON. Port : FROP. Le transport HTTP PEUT être protégé
par TLS~1.3 (recommandé) ; il n'existe pas d'exigence normative de confidentialité
des échanges ASAS entre GW.

\textbf{Certificat serveur (si TLS) :} Auto-signé ou CA locale -- aucune PKI
inter-opérateur.

\textbf{Adresse IP ASAS :} IP source post-NAT SIP (SD\_PUB\_IP, RT-0097 §6.2.2.4).

\textbf{Authentification normative :} Ancrage MCx + clés Ed25519. Tous les échanges
ASAS authentifiés par signature Ed25519. Pas de PKI inter-GW requise.

\textbf{Sérialisation :} UTF-8 + séparateurs \texttt{0x00} (cf.\ §\ref{sec:trust}).

\textbf{Timers :} recovery INTEGER(1..180)\,s (FROP) ; T\_KA INTEGER(1..10)\,s
(application profile) ; $N = 2$ (fixe).

\textbf{Option mTLS :} Accord bilatéral FSNNI OOB (FROP).

\paragraph{X.3 -- Initialisation (\texttt{POST /asas/appsessions})}

\begin{lstlisting}[language={}]
// Requete
{ "appSessionId_call": "uuid-v4", "svcSessionRef_1": "uuid-v4",
  "svcSessionId_1": "id-natif", "callingPublicKey": "base64(Ed25519)",
  "keepaliveReq": 2,
  "mspCapabilities": { "supportedPolicies": ["DUPLICATION","RECOVERY","LOADBALANCE"],
                       "maxPaths": 4 } }

// Reponse 200 OK
{ "appSessionId_term": "uuid-v4", "svcSessionRef_1": "uuid-v4",
  "terminatingPublicKey": "base64(Ed25519)", "keepaliveAccepted": 2,
  "mspPolicy": { "policy": "DUPLICATION|RECOVERY|LOADBALANCE|NONE", "activePaths": 2 },
  "advertisedAliases": [{ "svcSessionRef": "uuid-v4",
    "serviceType": "MCX|FUTURE_CRITICAL_SERVICE",
    "serviceDomain": "string", "serviceId": "string",
    "reachability": "CURRENT|ALTERNATIVE" }] }
\end{lstlisting}

\paragraph{X.4 -- Declaration / recovery}

\begin{lstlisting}[language={}]
{ "appSessionId_call": "uuid-v4", "appSessionId_term": "uuid-v4",
  "svcSessionRef": "uuid-v4", "serviceType": "MCX|FUTURE_CRITICAL_SERVICE",
  "svcSessionId": "id-natif", "svcSessRequestTime": "ISO8601",
  "svcSessAcceptTime": "ISO8601",
  "sessionToken": "base64(Sign(sk_call, quintuplet_0x00))" }
// Recovery : ajouter "replSvcSessionRef": "uuid-v4"
// Reponse : { "svcSessionRef": "uuid-v4", "ackToken": "base64(Sign(sk_term, ...))" }
\end{lstlisting}

\paragraph{X.5 -- Keepalive}
Envoye si silence $\geq$ T\_KA.
\begin{lstlisting}[language={}]
{ "appSessionId_call": "uuid-v4", "appSessionId_term": "uuid-v4",
  "svcSessionRefs": ["uuid-v4"], "timestamp": "ISO8601" }
\end{lstlisting}

\paragraph{X.6 -- Mise à jour alias (\texttt{PATCH .../aliases})}
\textbf{Immédiat} sur tout nouvel enregistrement. Sans configuration FROP.
\begin{lstlisting}[language={}]
{ "appSessionId_call": "uuid-v4", "appSessionId_term": "uuid-v4",
  "updatedAliases": [{ "svcSessionRef": "uuid-v4",
    "serviceType": "MCX|FUTURE_CRITICAL_SERVICE", "serviceDomain": "string",
    "serviceId": "string", "reachability": "CURRENT|ALTERNATIVE",
    "action": "ADD|REMOVE" }] }
\end{lstlisting}

\paragraph{X.7 -- Liberation}
\begin{lstlisting}[language={}]
{ "appSessionId_call": "uuid-v4", "appSessionId_term": "uuid-v4",
  "reason": "NORMAL|ERROR|TIMEOUT" }
\end{lstlisting}

\paragraph{X.8 -- Machines d'etat}
\begin{lstlisting}[language={}]
Session:
[CREATING] -- 200 OK --> [ESTABLISHED] -- all lost --> [DEGRADED]
[DEGRADED] -- >=1 ACTIVE --> [ESTABLISHED] | -- timer --> [RELEASING]
[ESTABLISHED|CREATING] -- DELETE --> [RELEASING] -- done --> [RELEASED]

Service session:
[CREATING] -- OK --> [ACTIVE] -- 2 KA fail --> [DEGRADED] -- recovery --> [ACTIVE]
[ACTIVE|DEGRADED] -- RELEASING --> [RELEASED]
\end{lstlisting}

\paragraph{X.9 -- Mapping erreurs \arr\ causes \OBAPP}

\begin{longtable}{p{5.5cm} p{8cm}}
  \toprule Cause ASAS & Cause \texttt{openSessionFinalAnswerNotif} \\ \midrule \endhead
  ASAS timeout & \texttt{MCX\_ENDPOINT\_NOT\_REACHABLE} \\
  ASAS 403 & \texttt{REMOTE\_ENDPOINT\_DECLINED} \\
  ASAS 404 & \texttt{TERMINATING\_APPLICATION\_ENDPOINT\_NOT\_REACHABLE} \\
  ASAS 401 & \texttt{TERMINATING\_APPLICATION\_ENDPOINT\_NOT\_ALLOWED} \\
  \bottomrule
\end{longtable}

% ── 5. TABLEAU RÉCAPITULATIF ────────────────────────────────────────────────
\section{Tableau récapitulatif des impacts}

\begin{longtable}{p{5cm} p{6.5cm} p{2cm}}
  \toprule Document & Nature de la modification & Specs UIC \\ \midrule \endhead
  \textbf{ETSI TS 103\,764} & Session + service session + AIR + MSP + ASAS +
    service identifier (§3.1) + ASM + \textbf{Annexe X} & \\
  \textbf{ETSI TS 103\,765-3} & Remplacement + Step 1
    (AIR \arr\ id de service + SD\_PUB\_IP) + MSP policy & inchangées \\
  \textbf{ETSI TS 103\,765-4} & Remplacement + Step 0 +
    alias list (déclenchement immédiat) + MSP & inchangées \\
  \textbf{ETSI TS 103\,765-5} & Remplacement + binding MSP +
    distinction Multipath/MSP & inchangées \\
  \textbf{UIC FFFIS-7950} & \textbf{Aucune modification} & \cmark \\
  \textbf{UIC FIS-7970} & \textbf{Aucune modification} & \cmark \\
  \textbf{UIC FRS FU-7120} & \textbf{Aucune modification} & \cmark \\
  \textbf{UIC AT-7800 SRS} & \textbf{Aucune modification} & \cmark \\
  \textbf{UIC TOBA FRS-7510} & \textbf{Aucune modification} & \cmark \\
  \textbf{3GPP TS 24.582/24.379} & \textbf{Aucune modification} & \\
  \bottomrule
\end{longtable}

% ── 6. POINTS OUVERTS ───────────────────────────────────────────────────────
\section{Points ouverts}

Tous les points ouverts sont résolus. Le document est prêt pour soumission ETSI.

\begin{enumerate}[leftmargin=2em]
  \item[\st{1.}] Adresse IP ASAS / NAT -- \textit{Résolu} (§\ref{sec:nat}).
  \item[\st{2.}] Authentification multi-hop FSNNI -- \textit{Résolu} (§\ref{sec:trust}).
  \item[\st{3.}] Sérialisation sessionToken -- \textit{Résolu}.
  \item[\st{4.}] Règles AIR -- \textit{Résolu} : \texttt{*.ertms}
    \arr\ \texttt{@current\_domain} ; sinon pass-through.
  \item[\st{5.}] Timer de recovery -- \textit{Résolu} :
    INTEGER(1..180)\,s ; $\leq 5$\,s pour ETCS/ATP.
  \item[\st{6.}] T\_KA / N -- \textit{Résolu} :
    INTEGER(1..10)\,s ; $N=2$ ; T\_KA $\leq$ timer/2.
  \item[\st{7.}] Déclenchement liste d'alias -- \textit{Résolu} : immédiat, sans FROP.
  \item[\st{8.}] Extensibilité -- \textit{Résolu} : champ \texttt{serviceType} ouvert.
\end{enumerate}

\end{document}